\section{实验结果与分析}
\label{sec:results}

本章主要是对前面提出的推荐模型做一个实验验证。我们在两个数据集上面跑了实验，分别是电商数据集和新闻数据集，数据集的基本情况见表~\ref{tab:data}。和之前的研究\cite{li2021,chen2022}一样，我们用留存率和点击率作为主要的评价指标。

\begin{table}[htbp]
  \centering
  \caption{数据集基本情况}
  \label{tab:data}
  \begin{tabular}{lrrr}
    \hline
    数据集 & 用户数 & 物品数 & 交互数 \\
    \hline
    电商 & 52,341 & 18,906 & 1,204,557 \\
    新闻 & 31,027 & 64,112 & 873,290 \\
    \hline
  \end{tabular}
\end{table}

模型的损失函数定义如式~\eqref{eq:loss}所示：
\begin{equation}
  \mathcal{L} = -\sum_{(u,i)} \log \sigma(\hat{y}_{ui}) + \lambda \|\Theta\|_2^2
  \label{eq:loss}
\end{equation}
其中$\lambda$是正则化系数，实验里我们把它设成了$10^{-4}$。

实验结果表明，在电商数据集上，我们的模型的七日保留率达到了41.7\%，比基线模型高了大概2.9个百分点；在新闻数据集上点击率也有一定的提升。另外我们还发现，用户的历史交互越多，模型的效果就越好，这说明历史交互数量的增加导致了推荐效果的提升\cite{wang2020}。

总的来说，本文的模型在留存率上的表现是比较好的，但是在冷启动用户上面效果还不是很理想，这个问题我们会在第~\ref{sec:future}节进一步讨论。
